\documentclass[11pt]{article}
\usepackage[T1]{fontenc}
\usepackage{textcomp}
\usepackage[margin=0.75in]{geometry}
\usepackage{amsmath,amssymb,amsthm}
\usepackage{array,booktabs}
\usepackage{hyperref}
\setlength{\parindent}{0pt}
\newtheorem{theorem}{Theorem}
\theoremstyle{definition}
\newtheorem{definition}{Definition}
\title{Flow Proof Artifact: Rat-add-cancel}
\author{Generated by \texttt{flow doc proof}}
\date{Flow Proof Book}
\begin{document}
\maketitle
Equal rationals can be subtracted from both sides.\\[0.5em]
\textbf{Source.} landau — *Foundations of Analysis*\\[0.5em]
\bigskip
\noindent{\large \textbf{Derived fact 1.} \textit{p + r = q + r implies p = q} \hfill Rat \cdot addition \cdot right cancellation holds}\par
\smallskip\noindent\textit{Coordinate: Rat \cdot addition \cdot right cancellation holds}\par
\smallskip\noindent\textit{Source: landau}\par
\smallskip\noindent\textit{Built on: parentheses do not matter, for addition on Rat, zero is the right identity, for addition on Rat}\par
\medskip
\noindent\textbf{Goal.} p + r = q + r implies p = q\par
\noindent$\displaystyle \forall p \in Rat \forall q \in Rat \forall r \in Rat\quad p = q$\par
\medskip
\noindent
\begin{tabular}{@{} >{\raggedright\arraybackslash}p{0.06\textwidth} >{\raggedright\arraybackslash}p{0.47\textwidth} >{\raggedleft\arraybackslash}p{0.41\textwidth} @{}}
\textbf{\#} & \textbf{Proof} & \textbf{Mathematics} \\
\midrule
\textbf{1.} & We prove that right cancellation holds for addition on Rat. & \\[0.45em]
\textbf{2.} & We split into exhaustive cases — the claim must hold in each one. & \\[0.45em]
\textbf{3.} & Case 1 (see step 2): suppose p + r  equals  q + r. & \\[0.45em]
\textbf{4.} & We invoke the derived fact governing addition on Rat: parentheses do not matter, for addition on Rat (instantiated for p, r, -r). & \\[0.45em]
\textbf{5.} & We invoke the definitional clause governing addition on Rat: zero is the right identity, for addition on Rat (instantiated for p). & \\[0.45em]
\textbf{6.} & From step 3, step 4, and step 5, this implies p equals q. Together with the other cases (step 3), the goal is discharged. Hence proven. & $\displaystyle p = q$ \\[0.45em]
\bottomrule
\end{tabular}
\label{thm:Rat-addition-right_cancellation_holds}
\par\medskip\hrule
\end{document}
